\chapter{Utilizador e task model principal}
\label{cap:utilizador}

O utilizador principal é o supervisor, porque é ele quem cria os percursos, acompanha os estudantes e alimenta, com as suas avaliações, o progresso e as badges que os estudantes veem. Sem a sua atividade regular, a área do supervisionado fica vazia.

\section{Utilizador principal}

A Tabela~\ref{tab:utilizador} resume o perfil-tipo do supervisor.

\begin{table}[H]
    \centering
    \small
    \begin{tabularx}{\linewidth}{lX}
        \toprule
        \textbf{Aspeto} & \textbf{Descrição} \\
        \midrule
        Quem é & Enfermeiro ou docente que orienta estudantes em contexto de bloco operatório, normalmente vários em simultâneo e em percursos diferentes. \\
        \midrule
        Características & Domina a matéria clínica e os critérios de avaliação; tem competência digital média; tem pouco tempo entre turnos; usa o computador do serviço e, por vezes, um tablet ou telemóvel. \\
        \midrule
        Necessidades & Montar planos sem os reescrever de raiz; ver de relance em que ponto está cada estudante; registar avaliações logo após a sessão prática, de forma rápida e sem ambiguidade. \\
        \midrule
        Objetivos & Garantir que cada estudante passa por todas as fases de aprendizagem dentro dos prazos; ter um registo das avaliações que possa justificar; manter os estudantes motivados ao longo do percurso. \\
        \midrule
        Frustrações atuais & Folhas de cálculo diferentes para cada grupo; progresso calculado à mão; avaliações em papel que se perdem ou ficam por registar. \\
        \bottomrule
    \end{tabularx}
    \caption{Perfil do utilizador principal: o supervisor.}
    \label{tab:utilizador}
\end{table}

O supervisionado é um utilizador secundário: consulta o seu plano e o seu progresso, mas não altera dados. O administrador usa a aplicação com pouca frequência, sobretudo no início de cada ano letivo.

\section{Tarefa principal}

A tarefa principal é \textbf{avaliar um supervisionado após uma sessão prática}. É a tarefa mais frequente do supervisor e é ela que faz avançar o progresso do estudante e que desencadeia a atribuição de badges (Figura~\ref{fig:bpmn-avaliacao}).

A Figura~\ref{fig:task-model} decompõe a tarefa numa análise hierárquica de tarefas (\gls{hta})~\citep{annett-hta}. Cada nível divide a tarefa de cima em passos mais concretos, e os planos indicam a ordem em que esses passos são feitos.

\begin{figure}[H]
    \centering
    \resizebox{\linewidth}{!}{% Task model (HTA): avaliar um supervisionado após uma sessão prática.
\begin{tikzpicture}[sub/.style={tarefahta, fill=white, text width=2.2cm, minimum height=0.9cm}]
  % Tarefa principal
  \node[tarefaraiz] (r) at (5.85,2.3) {\textbf{0.} Avaliar um supervisionado após uma sessão prática};

  % Subtarefas
  \node[tarefahta] (n1) at (0,0)    {\textbf{1.} Entrar na aplicação};
  \node[tarefahta] (n2) at (3.9,0)  {\textbf{2.} Abrir a ficha do supervisionado};
  \node[tarefahta] (n3) at (7.8,0)  {\textbf{3.} Classificar as componentes};
  \node[tarefahta] (n4) at (11.7,0) {\textbf{4.} Guardar e confirmar};

  \draw[gInk, line width=0.6pt] (r.south) -- (5.85,1.15);
  \draw[gInk, line width=0.6pt] (0,1.15) -- (11.7,1.15);
  \foreach \nd/\x in {n1/0, n2/3.9, n3/7.8, n4/11.7}
    \draw[gInk, line width=0.6pt] (\x,1.15) -- (\nd.north);

  % Ações
  \node[sub] (c11) at (0.6,-1.5)  {\textbf{1.1} Inserir email e palavra-passe};
  \node[sub] (c12) at (0.6,-2.8)  {\textbf{1.2} Carregar em \emph{Login}};

  \node[sub] (c21) at (4.5,-1.5)  {\textbf{2.1} No painel, abrir o percurso};
  \node[sub] (c22) at (4.5,-2.8)  {\textbf{2.2} Na lista de inscritos, abrir o aluno};

  \node[sub] (c31) at (8.4,-1.5)  {\textbf{3.1} Abrir um módulo};
  \node[sub] (c32) at (8.4,-2.8)  {\textbf{3.2} Escolher a fase};
  \node[sub] (c33) at (8.4,-4.1)  {\textbf{3.3} Escolher o estado de cada componente};

  \node[sub] (c41) at (12.3,-1.5) {\textbf{4.1} Carregar em \emph{Guardar Avaliações}};
  \node[sub] (c42) at (12.3,-2.8) {\textbf{4.2} Confirmar o progresso na lista de inscritos};

  \begin{scope}[gInk, line width=0.5pt]
    \draw ($(n1.south)+(-1.0,0)$) |- (c11.west);
    \draw ($(n1.south)+(-1.0,0)$) |- (c12.west);
    \draw ($(n2.south)+(-1.0,0)$) |- (c21.west);
    \draw ($(n2.south)+(-1.0,0)$) |- (c22.west);
    \draw ($(n3.south)+(-1.0,0)$) |- (c31.west);
    \draw ($(n3.south)+(-1.0,0)$) |- (c32.west);
    \draw ($(n3.south)+(-1.0,0)$) |- (c33.west);
    \draw ($(n4.south)+(-1.0,0)$) |- (c41.west);
    \draw ($(n4.south)+(-1.0,0)$) |- (c42.west);
  \end{scope}

  % Planos
  \node[plano, anchor=north west, text width=14.5cm] at (-1.35,-4.9) {%
    \textbf{Plano 0:} fazer 1--2--3--4; o passo 1 só é necessário se não houver sessão iniciada.\\
    \textbf{Plano 3:} para cada módulo, fazer 3.1; depois, para cada fase prática, fazer 3.2--3.3.\\
    \textbf{Plano 4:} fazer 4.1--4.2; se o progresso não mudou, voltar a 3.};
\end{tikzpicture}
}
    \caption{Task model da tarefa principal, em HTA.}
    \label{fig:task-model}
\end{figure}

A sequência de ações, a partir do ecrã de entrada, é a seguinte:

\begin{enumerate}
    \item Na página de entrada, o supervisor insere o email e a palavra-passe e carrega em \emph{Login}; a aplicação abre o painel do supervisor (\texttt{/dashboard}).
    \item No painel, carrega no cartão do percurso; abre-se a gestão do percurso, com a lista de supervisionados inscritos (\texttt{/pathway/\{id\}/students}).
    \item Na lista, carrega no estudante; abre-se a ficha de avaliação (\texttt{/evaluate/\{id\}}).
    \item Na ficha, abre um módulo, escolhe a fase prática e, para cada componente, carrega num dos estados: \emph{Inconsistente}, \emph{Abaixo da média}, \emph{Acima da média} ou \emph{Consistente}.
    \item Repete o passo anterior para cada fase e cada módulo que avaliou na sessão.
    \item Carrega em \emph{Guardar Avaliações}. O sistema grava as classificações, recalcula o progresso, atribui as badges devidas e notifica o estudante.
    \item Volta à lista de inscritos e confirma o novo valor de \emph{Progresso no Percurso}.
\end{enumerate}

Sem contar com a autenticação, a tarefa exige três ecrãs e duas navegações até à ficha de avaliação. O resto do trabalho acontece num único ecrã, como pede o requisito RNF01.
